\documentclass[10pt,a4paper,sans]{moderncv}
\moderncvstyle{classic}
\moderncvcolor{blue}
\usepackage[utf8]{inputenc}
\usepackage[scale=0.82]{geometry}

% moderncv compose \today en anglais par défaut ; format français sans charger babel.
\renewcommand{\today}{\number\day\space\ifcase\month\or janvier\or février\or mars\or avril\or
mai\or juin\or juillet\or août\or septembre\or octobre\or novembre\or décembre\fi\space\number\year}

\name{Alexandro}{Disla}
\title{Lettre de Motivation}
\address{Rue Saint-Louis Jeanty \#14, Route de Frère}{Port-au-Prince, Haïti}{}
\phone[mobile]{(+509) 4148-3700}
\email{alexandrodisla@hotmail.com}
\extrainfo{GitHub: \url{https://github.com/AD0791}}

\begin{document}

\recipient{Équipe de Recrutement --- Catholic Relief Services Haïti}{Port-au-Prince, Haïti\\Poste : Project Officer MEAL --- Projet RAPID}
\date{\today}
\opening{Madame, Monsieur,}
\closing{Je vous prie d'agréer, Madame, Monsieur, l'expression de mes salutations distinguées,}
\enclosure[Ci-joint]{Curriculum Vitae détaillé, copies de diplômes}

\makelettertitle

Je vous adresse ma candidature au poste de Project Officer MEAL du projet RAPID. Par souci de transparence : j'ai également adressé ma candidature au poste de Senior Project Officer MEAL du même projet. Les deux m'intéressent réellement, et je serais heureux de contribuer à RAPID à l'un ou l'autre niveau.

Ce poste décrit mon travail quotidien depuis cinq ans. À la Caris Foundation International, j'ai conduit le suivi d'un projet de santé et de VIH sur des cliniques multi-sites : conception et standardisation des formulaires, encadrement technique des agents de collecte sur l'ensemble des sites, vérification des données remontées, calcul des indicateurs et soumission du rapportage. J'ai construit les instruments de collecte sur CommCare, KoboToolbox, ODK et mWater --- avec logiques de saut et contraintes de validation, pour que l'erreur soit arrêtée à la saisie plutôt que corrigée trois semaines plus tard --- et j'ai formé les agents de collecte et les équipes partenaires à leur usage, guides techniques à l'appui, de manière qu'elles continuent sans moi. Sur la qualité des données, j'ai établi et supervisé des protocoles de DQA en remontant les chiffres rapportés jusqu'aux registres sources, et en suivant les actions correctives jusqu'à leur clôture. Enfin, je produisais et soumettais les indicateurs MER destinés à l'USAID via DATIM : le rythme, les contrôles de cohérence et les échéances d'un bailleur du gouvernement américain me sont familiers, ce qui devrait servir un projet financé par le Département d'État.

Vous verrez sur mon CV que mon expérience dépasse les deux ans demandés, et je préfère dire pourquoi je postule à ce poste plutôt que de laisser la question en suspens. J'exerce aussi comme ingénieur de données : j'administre les bases et je construis les tableaux de bord Power BI dont un système MEAL dépend, au lieu de seulement les spécifier. Ce que décrit votre fiche --- vérification quotidienne de la qualité, nettoyage des bases, appui aux équipes sur les outils numériques --- est précisément l'endroit où cette double compétence sert le plus, et c'est ce travail-là que je veux continuer à faire plutôt que de m'en éloigner.

Deux réserves, que je préfère poser franchement. Ma pratique du suivi s'est faite en santé et VIH, en eau et assainissement et en éducation, dans des projets de développement plutôt qu'en réponse d'urgence : je n'ai pas conduit de PDM, de suivi de distributions ni de vérification post-enregistrement, exercices propres à l'assistance alimentaire, même si l'enquête, l'échantillonnage et la vérification qui les composent me sont familiers. De même, j'ai travaillé dans le respect des principes de redevabilité sans avoir moi-même opéré un mécanisme de gestion des plaintes ; c'est un volet que j'aborderais en apprenant vos procédures, pas en prétendant les connaître.

Ce projet m'attire parce que son approche convergente --- les mêmes ménages suivis à travers plusieurs services essentiels --- pose exactement le problème de données que je sais traiter : identifier un ménage de façon fiable d'un secteur à l'autre, sans le compter deux fois ni le perdre en chemin. Je reste à votre entière disposition pour un entretien.

\makeletterclosing

\end{document}
